\section{Introduction}\label{sec:introduction}

Online social networks (OSNs), such as Twitter/X, Facebook, and Weibo, have become important infrastructures for information sharing and public discussion. Social bots are automated or semi-automated accounts that mimic human behavior and interact with users at scale~\cite{ferrara2016rise}. Malicious bots can spread misinformation, amplify coordinated opinions, invade privacy, and manipulate online topics, making reliable social bot detection an important task~\cite{feng2021twibot20,twibot22benchmark}.

Social bot detection has shifted from handcrafted feature engineering to representation learning over social graphs. Early detectors used metadata, activity patterns, network statistics, and tweet content with conventional classifiers~\cite{davis2016botornot,varol2017online,bot_hunter2018,friendbot2020}, but observable profile and content signals are vulnerable to adaptive manipulation and content imitation~\cite{cresci2017paradigm}. Recent methods therefore exploit account attributes, text, and social relations: heterogeneous graph detectors model relation-specific interactions~\cite{feng2021botrgcn,feng2022rgt}; text--graph and multimodal methods fuse textual and structural evidence~\cite{lei2023bic,liu2023botmoe,zhou2024lgb}; reliability-aware approaches reduce noisy or heterophilous edges~\cite{qiao2025bece,lin2025botbr}; and higher-order methods capture group-wise relations beyond pairwise links~\cite{gao2025higherorder,yang2024sebot}. These studies strengthen detector representations, but they leave a complementary question underexplored: how to estimate node-level prediction reliability and determine when a detector's prediction should be trusted.

To address this gap, we propose BotRHG, a reliability-guided hypergraph learning framework for social bot detection. A base detector provides initial node-level decisions, a local reliability estimate identifies low-reliability accounts, and a hypergraph supplies higher-order evidence only for routed accounts. For routed accounts, selective residual correction fuses the base representation with hypergraph evidence from semantically or structurally related users; for non-routed accounts, the base prediction is preserved.

Our experiments on TwiBot-20 and TwiBot-22 demonstrate that the proposed method achieves state-of-the-art performance, outperforming eight competitive social bot detection baselines. Extensive analyses further validate the effectiveness of reliability-guided routing and selective residual correction, and show that selective higher-order evidence can improve the detection of low-reliability accounts.

The main contributions of this work are summarized as follows:
\begin{enumerate}
\item We propose BotRHG, a reliability-guided hypergraph learning framework for social bot detection, using node-level prediction reliability to guide selective higher-order learning.
\item We develop a reliability-guided correction module that converts local reliability estimates into account-level routing decisions and performs selective residual fusion over support hyperedges, enabling targeted higher-order correction while preserving reliable base predictions.
\item Our experiments show that the proposed method achieves state-of-the-art performance on two widely used Twitter bot benchmarks. Extensive analyses further validate the effectiveness and robustness of our approach.
\end{enumerate}
